%% ════════════════════════════════════════════════════════════════════════════
\section{Methodology}
\label{sec:method}

\subsection{Experimental Setups}
To decouple the components of predictive performance and quantify artificial inflation, we define three nested experimental configurations executed over municipal series (with Setups~A and B restricted to valid-PM$_{2.5}$ days, and Setup~C evaluated over the full valid-AQI record as well as a like-for-like 230-day sensitivity check on Setup~B's window):

\begin{description}
  \item[Setup A (Random Split, Concurrent Inputs):] Features consist of concurrent same-day criteria pollutant concentrations $\{PM_{2.5,t}, PM_{10,t}, NO_{2,t}, SO_{2,t}, CO_t, O_{3,t}, NH_{3,t}\}$ and calendar harmonic encodings. Training and test samples are randomly partitioned via an 80/20 ratio across 10 independent random seeds. This mirrors the prevailing setup in published literature and captures the combined effect of concurrent input leakage and autocorrelation leakage.
  \item[Setup B (Chronological Split, Concurrent Inputs):] Utilises the identical concurrent feature matrix as Setup~A, but enforces a strict chronological 80/20 partition: the initial 80\% of calendar days form the training set, and the final 20\% form the test set. By holding the feature set constant while restoring temporal ordering, Setup~B reflects the combined effect of eliminating random-split autocorrelation leakage and any temporal distribution shift between the training and test windows.
  \item[Setup C (Rolling-Origin, Strictly Lagged Inputs):] Evaluates true advance forecasting ($h=1$ and $h=3$ days ahead). Predictor variables $X_t$ are restricted strictly to information recorded at or prior to forecast origin $t$. Models are evaluated over the final 20\% temporal window using an expanding training origin refit every 30 calendar days.
\end{description}

\subsection{Feature Engineering and Runtime Leakage Guards}
For Setup~C, a unified 49-column feature schema was engineered from strictly historical information:
\begin{itemize}
  \item \textbf{Autoregressive Target Lags:} $\{y_t, y_{t-1}, y_{t-2}, y_{t-6}, y_{t-13}\}$ (indexed strictly relative to forecast origin $t$).
  \item \textbf{Lagged Criteria Pollutants:} 1, 2, 3, and 7-day lagged concentrations of the seven criteria pollutants (28 columns). In the four cities without Kaggle pollutant records (Gandhinagar, Ankleshwar, Surat, Vapi), these 28 columns are entirely missing and imputed/zero-filled, leaving 21 active AQI-lag, rolling-statistic, and calendar features. Moreover, because Kaggle pollutant records end in 2020, these 28 pollutant-lag columns are also missing and imputed during the 2022--2023 Setup~C test windows for Ahmedabad and Chennai.
  \item \textbf{Rolling Historical Statistics:} 3, 7, 14, and 30-day backward-looking rolling means and standard deviations of AQI (8 columns).
  \item \textbf{Calendar Harmonic Encodings:} Cyclical sine and cosine transformations of day-of-week, month, and day-of-year corresponding to the target forecast date $t+h$, plus day-of-week/weekend indicators (8 columns).
\end{itemize}

To enforce pipeline integrity, the codebase embeds three runtime automated leakage guards:
\begin{enumerate}
  \item \emph{Temporal Separation Guard:} Formally asserts $\max(\text{train\_dates}) < \min(\text{test\_dates})$ across all evaluation folds.
  \item \emph{Planted-Leak Detector:} Computes pairwise correlations between all candidate features and out-of-sample targets, raising an exception if any feature achieves $|r| \ge 0.9999$.
  \item \emph{Shuffled-Target Falsification Guard:} Verifies that randomly permuting the training target vector completely destroys predictive performance ($R^2 \le 0.05$).
\end{enumerate}

\subsection{Candidate Forecasters}
We benchmark six primary forecasters spanning classical naive baselines, linear regularised estimators, and nonlinear ensemble architectures:
\begin{enumerate}
  \item \textbf{Persistence Baseline:} Assumes zero atmospheric change: $\hat{y}_{t+h} = y_t$.
  \item \textbf{Seasonal Naive Baseline:} Projects the value from the identical day of the preceding week: $\hat{y}_{t+h} = y_{t-7+h}$.
  \item \textbf{7-Day Moving Average:} Projects the unweighted mean of the preceding week: $\hat{y}_{t+h} = \frac{1}{7}\sum_{k=0}^{6} y_{t-k}$.
  \item \textbf{Ridge Regression:} $L_2$-regularised linear model ($\alpha = 1.0$) trained on standard-scaled lagged predictors.
  \item \textbf{Random Forest:} Ensemble of 100 decision trees, $\text{max\_features} = \sqrt{p}$, $\text{min\_samples\_leaf} = 5$.
  \item \textbf{XGBoost:} Gradient boosted trees (100 estimators, learning rate $\eta = 0.1$, max depth $= 4$, subsample ratio $= 0.8$).
\end{enumerate}
Primary hyperparameters were established based on environmental machine learning conventions and held fixed across all cities to eliminate post-hoc tuning leakage; to test sensitivity to fixed defaults, we also report a 3-fold \texttt{TimeSeriesSplit} tuned check on Ahmedabad in Section~\ref{sec:dm_synthesis}. In addition, to ensure ML models are compared against standard statistical benchmarks at $h=1$, we evaluate four supplementary classical baselines refit at the identical 30-day expanding cadence: (7)~\textbf{Monthly Climatology} (historical mean for the target calendar month computed on the expanding training fold), (8)~\textbf{Climatology+Persistence} (an unweighted 50/50 blend), (9)~\textbf{ARIMA(1,0,0)} (an AR(1) model fit on the expanding training series and rolled forward one step from $y_t$), and (10)~\textbf{ETS (Simple Exponential Smoothing)} with estimated smoothing level $\alpha$.

\subsection{Evaluation Metrics and Hypothesis Testing}
Primary predictive accuracy is measured using Mean Absolute Error (MAE), supplemented by Root Mean Squared Error (RMSE) and $R^2$. Band Accuracy (Band Acc) measures the fraction of predictions that fall within the same CPCB categorical AQI band as the observed value (Good, Satisfactory, Moderate, Poor, Very Poor, or Severe). To account for serial correlation in error residuals, 95\% confidence intervals for MAE are estimated via the moving block bootstrap ($B = 1{,}000$ resamples, block length $b = 7$ days matching the empirical weekly cycle).

The final 20\% of each city's valid-day record forms the test window. This corresponds to approximately 65 test points for Surat ($n = 327$ valid days) and roughly one season ($\sim$340 days) for Gandhinagar; smaller test windows reduce Diebold-Mariano power and should be interpreted with appropriate caution.

To test whether model improvements over Persistence are statistically significant, we deploy the two-tailed \textbf{Diebold-Mariano (DM) test} \cite{diebold1995}. Let $e_{1,t} = y_{t+h} - \hat{y}_{1,t+h}$ denote the error of the Persistence baseline, and let $e_{2,t} = y_{t+h} - \hat{y}_{2,t+h}$ denote the error of the candidate model. The loss differential under absolute loss is:
\begin{equation}
  d_t = |e_{1,t}| - |e_{2,t}|
  \label{eq:dm_diff}
\end{equation}
The DM test statistic is:
\begin{equation}
  \text{DM} = \frac{\bar{d}}{\sqrt{\hat{V}(\bar{d})}}
  \label{eq:dm_stat}
\end{equation}
where $\bar{d} = \frac{1}{T}\sum_{t=1}^T d_t$, and $\hat{V}(\bar{d})$ is the rectangular autocovariance-truncated variance estimator up to lag $k = h - 1$ with the Harvey--Leybourne--Newbold (HLN) small-sample correction. At $h=1$, $k=0$ and no autocorrelation correction is applied; the test reduces to an HLN small-sample corrected $t$-test on the loss differentials. At $h=3$, $k=2$ lags are included to account for overlapping multi-step forecast dependence. As a sensitivity check at $h=1$, we also evaluate $k=3$ lags. Under the null hypothesis of equal forecast accuracy ($H_0: E[d_t] = 0$), the HLN-corrected statistic is referenced against a Student-$t$ distribution with $T - 1$ degrees of freedom. Two-tailed $p$-values below $0.05$ indicate nominally significant differences; for the primary family of 60 candidate-vs-Persistence comparisons (6 cities $\times$ 5 primary candidates $\times$ 2 horizons), the Bonferroni-corrected threshold is $\alpha = 0.05 / 60 = 0.00083$.

